\section{Протоколна основа}
\label{sec:basis}

\subsection{Записният слой на TLS}

TLS пренася всичко, което изпраща, в записи. Всеки запис започва с петоктетна заглавна
част, а след нея идва съдържанието. В TLS 1.3 структурата се нарича \code{TLSPlaintext}
и е определена в раздел 5.1 на RFC 8446 \cite{rfc8446}. В TLS 1.2 същата структура е
описана в раздел 6.2.1 на RFC 5246 \cite{rfc5246}. Полетата са следните.

\begin{table}[H]
\centering
\caption{Заглавна част на TLS запис}
\label{tab:record}
\begin{tabularx}{\textwidth}{@{}lcY@{}}
\toprule
Поле & Октети & Съдържание \\
\midrule
\code{type} (ContentType) & 1 & вид на съдържанието: change\_cipher\_spec (20), alert (21),
  handshake (22) или application\_data (23) \\
\code{legacy\_record\_version} & 2 & в TLS 1.3 трябва да е 0x0303; в първия ClientHello
  може да е и 0x0301 заради съвместимост \\
\code{length} & 2 & дължина на съдържанието в октети; не може да надхвърля $2^{14}$ \\
\code{fragment} & \code{length} & самото съдържание \\
\bottomrule
\end{tabularx}
\end{table}

Първото поле е определящо за цялата работа. Стойностите му са малко и са фиксирани:
промяна на шифъра 20 (0x14), предупреждение 21 (0x15), ръкостискане 22 (0x16) и
приложни данни 23 (0x17). В TLS клиентът говори пръв. Първото, което той изпраща, е
съобщението ClientHello, а то е съобщение за ръкостискане. Следователно първият октет на
всяка TLS връзка е 0x16. Шестият октет пък е типът на съобщението за ръкостискане, който
за ClientHello е 1.

RFC 8446 изисква полето за версия да се игнорира и задава горна граница на дължината:
запис, по-дълъг от $2^{14}$ октета, трябва да прекъсне връзката с предупреждение
record\_overflow. Тези проверки са работа на TLS библиотеката. Разпознаването в
\code{octet} не ги дублира и в раздел~\ref{sec:security} се обяснява защо.

\subsection{Ръкостискането в TLS 1.3}

Пълното ръкостискане в TLS 1.3 изисква едно отиване и връщане. Клиентът изпраща
ClientHello, в което предлага поддържаните версии, параметри за обмен на ключове, по
желание името на сървъра (SNI) и по желание списък с приложни протоколи (ALPN). Сървърът
отговаря с ServerHello и от този момент нататък всичко е шифровано: допълнителните
разширения (включително избрания приложен протокол), сертификатът, подписът върху
ръкостискането и съобщението Finished. Клиентът проверява сертификата и подписа и
изпраща своето Finished \cite{rfc8446}.

Важно за тази работа е, че всичко преди ServerHello пътува некриптирано. Ако сървърът
откаже връзката още при обработката на ClientHello, предупреждението, с което отказва,
също е некриптирано и се вижда в мрежата. Това е използвано в теста на ALPN в
раздел~\ref{sec:verification}: там отказът се проверява по точните октети на записа.

\subsection{Предупрежденията}

Предупреждението (alert) е запис с ContentType 21 и съдържание от два октета: ниво и
описание \cite{rfc8446}. Ниво 2 означава фатално предупреждение, след което връзката се
затваря. Описанията, които се срещат в тази работа, са три. no\_application\_protocol
(120) означава, че сървърът не поддържа нито един от предложените приложни протоколи.
record\_overflow (22) означава твърде дълъг запис. unexpected\_message (10) означава
съобщение, което не е на мястото си.

\subsection{Договаряне на приложния протокол}

TLS защитава байтове, но не казва какво означават те. Над една и съща TLS връзка може да
се говори HTTP/1.1, HTTP/2 или друг протокол. Разширението ALPN решава този въпрос в
самото ръкостискане \cite{rfc7301}. Клиентът изпраща списък от имена. Всяко име е
предшествано от един октет с дължината му и празни имена не са позволени (раздел 3.1 на
RFC 7301). Сървърът избира едно име от списъка. Ако не поддържа нито едно, той отговаря с
фаталното предупреждение no\_application\_protocol (раздел 3.2).

Препоръките на IETF за сигурна употреба на TLS изискват от сървърите строго да спазват
това правило \cite[раздел 3.8]{rfc9325}. Причината е защитата от атаки между протоколи,
при които съобщение, предназначено за един протокол, бива прието от друг. Ако сървърът
приема връзката въпреки липсата на общ протокол, той дава възможност TLS сесия,
предназначена за друга услуга, да бъде насочена към него.

\subsection{Сертификати и верига на доверие}

Сървърът доказва самоличността си със сертификат по стандарта X.509 \cite{rfc5280}.
Сертификатът свързва публичен ключ с име и е подписан от удостоверяващ орган. Клиентът
има списък с органи, на които вярва. Той проверява, че веригата от сертификата на сървъра
до някой от тези органи е валидна и че името в сертификата съвпада с името, към което
клиентът се свързва.

В сертификата има разширения, които ограничават употребата му. Основните ограничения
(basicConstraints) казват дали сертификатът е на удостоверяващ орган. Употребата на ключа
(keyUsage) и разширената употреба (extendedKeyUsage) казват за какво служи ключът.
Алтернативното име на субекта (subjectAltName) съдържа имената, за които сертификатът
важи \cite{rfc5280}. В раздел~\ref{sec:mechanism} се описва как \code{octet} създава за
тестовете собствен орган и сертификат на сървъра с точно тези разширения.

\subsection{Началото на HTTP заявката}

HTTP/1.1 заявката започва с ред, съставен от метод, интервал, цел на заявката, интервал
и версия на протокола \cite[раздел 3]{rfc9112}. Методът е последователност от знаци,
която е чувствителна към регистъра. Всички стандартни методи са записани с главни букви
по споразумение \cite[раздел 9.1]{rfc9110}: GET, HEAD, POST, PUT, DELETE, CONNECT,
OPTIONS, TRACE. HTTP/2 без TLS може да започне директно с фиксирана начална
последователност, която започва с \code{PRI * HTTP/2.0} \cite[раздел 3.4]{rfc9113}.
Тя също започва с главна буква.

\subsection{Защо един октет е достатъчен}

От двата предходни подраздела следва, че началото на връзката попада в едно от две
непресичащи се множества. TLS връзката започва с 0x16. Некриптираната HTTP заявка
започва с главна латинска буква, тоест с октет между 0x41 и 0x5A. Октетът 0x16 е
управляващ знак в ASCII и не е буква. Таблица~\ref{tab:octets} показва как
\code{octet} разпределя всичките 256 стойности.

\begin{table}[H]
\centering
\caption{Как \code{octet} разпределя първия октет}
\label{tab:octets}
\begin{tabularx}{\textwidth}{@{}lYl@{}}
\toprule
Първи октет & Значение & Решение \\
\midrule
от 0x00 до 0x13 & управляващи знаци & отказ \\
0x14, 0x15, 0x17 & TLS запис, който не може да отвори връзка & отказ \\
0x16 & TLS запис с ръкостискане, тоест ClientHello & TLS \\
от 0x18 до 0x40 & управляващи знаци, интервал, цифри, препинателни знаци & отказ \\
от 0x41 до 0x5A & главни латински букви, начало на HTTP метод & HTTP \\
от 0x5B до 0x7F & малки букви и други знаци & отказ \\
от 0x80 до 0xFF & извън ASCII, включително старият формат на SSL 2.0 & отказ \\
\bottomrule
\end{tabularx}
\end{table}

Две стойности заслужават отделна бележка. Записите с ContentType 20, 21 и 23 са валидни
TLS записи, но никой от тях не може да е първият запис на нова връзка. Ако връзка започне
с тях, другата страна не започва TLS сесия, затова \code{octet} не я подава на TLS
библиотеката. Старият формат на ClientHello от SSL 2.0 започва с октет със зададен
най-старши бит. SSL 2.0 е забранен \cite{rfc6176} и \code{octet} не го разпознава като TLS.

Малките букви също водят до отказ. Това не е пропуск: методът е чувствителен към
регистъра и \code{get} не е методът GET. Сървър, който приема \code{get}, приема нещо,
което не е HTTP заявка.
